\documentclass[11pt,letterpaper]{article}

% --- Packages ---
\usepackage[margin=1in]{geometry}
\usepackage{graphicx}
\usepackage{booktabs}
\usepackage{longtable}
\usepackage{tabularx}
\usepackage{enumitem}
\usepackage{hyperref}
\usepackage{xcolor}
\usepackage{fancyhdr}
\usepackage{listings}
\usepackage{titlesec}
\usepackage[T1]{fontenc}
\usepackage{lmodern}
\usepackage{parskip}
\usepackage{array}
\usepackage{microtype}
\usepackage{tcolorbox}
\usepackage{amsmath}

% --- Colours ---
\definecolor{codebg}{gray}{0.95}
\definecolor{codeframe}{gray}{0.75}
\definecolor{linkblue}{HTML}{0066CC}
\definecolor{cautionred}{HTML}{CC0000}
\definecolor{notegray}{HTML}{555555}
\definecolor{guicolor}{HTML}{2E7D32}
\definecolor{tipgreen}{HTML}{1B5E20}
\definecolor{tipbg}{HTML}{E8F5E9}
\definecolor{warnbg}{HTML}{FFF3E0}
\definecolor{warnorange}{HTML}{E65100}
\definecolor{infobg}{HTML}{E3F2FD}
\definecolor{infoblue}{HTML}{0D47A1}
\definecolor{claudebg}{HTML}{F3E5F5}
\definecolor{claudepurple}{HTML}{6A1B9A}

% --- Hyperlink styling ---
\hypersetup{
    colorlinks=true,
    linkcolor=linkblue,
    urlcolor=linkblue,
    citecolor=linkblue,
    pdftitle={X-CAVATE CLI & Python API Reference},
    pdfauthor={X-CAVATE Team},
    pdfsubject={CLI and Python API Reference for X-CAVATE},
}

% --- Code listing style ---
\lstset{
    backgroundcolor=\color{codebg},
    frame=single,
    rulecolor=\color{codeframe},
    basicstyle=\ttfamily\small,
    breaklines=true,
    breakatwhitespace=false,
    columns=fullflexible,
    keepspaces=true,
    showstringspaces=false,
    xleftmargin=4pt,
    xrightmargin=4pt,
    aboveskip=8pt,
    belowskip=8pt,
}

\lstdefinestyle{python}{
    language=Python,
    basicstyle=\ttfamily\small,
    keywordstyle=\color{linkblue}\bfseries,
    stringstyle=\color{guicolor},
    commentstyle=\color{notegray}\itshape,
    backgroundcolor=\color{codebg},
    frame=single,
    rulecolor=\color{codeframe},
    breaklines=true,
    columns=fullflexible,
    keepspaces=true,
    showstringspaces=false,
    xleftmargin=4pt,
    xrightmargin=4pt,
    aboveskip=8pt,
    belowskip=8pt,
}

% --- Header / Footer ---
\pagestyle{fancy}
\fancyhf{}
\fancyhead[L]{\textbf{X-CAVATE CLI \& Python API Reference}}
\fancyhead[R]{\thepage}
\renewcommand{\headrulewidth}{0.4pt}

% --- Custom commands ---
\newcommand{\caution}[1]{\par\textcolor{cautionred}{\textbf{CAUTION:} #1}\par}
\newcommand{\note}[1]{\par\textcolor{notegray}{\textbf{Note:} #1}\par}
\newcommand{\code}[1]{\texttt{#1}}
\newcommand{\gui}[1]{\textcolor{guicolor}{\textsf{\textbf{#1}}}}
\newcommand{\btn}[1]{\fbox{\textsf{#1}}}
\newcolumntype{L}[1]{>{\raggedright\arraybackslash}p{#1}}
\newcolumntype{C}[1]{>{\centering\arraybackslash}p{#1}}
\tolerance=400
\emergencystretch=2em

% --- Callout boxes ---
\newtcolorbox{tipbox}{
    colback=tipbg,
    colframe=tipgreen,
    coltitle=white,
    fonttitle=\bfseries,
    title=Tip,
    boxrule=0.5pt,
    arc=2pt,
    left=6pt, right=6pt, top=4pt, bottom=4pt,
}

\newtcolorbox{warnbox}{
    colback=warnbg,
    colframe=warnorange,
    coltitle=white,
    fonttitle=\bfseries,
    title=Warning,
    boxrule=0.5pt,
    arc=2pt,
    left=6pt, right=6pt, top=4pt, bottom=4pt,
}

\newtcolorbox{infobox}{
    colback=infobg,
    colframe=infoblue,
    coltitle=white,
    fonttitle=\bfseries,
    title=Info,
    boxrule=0.5pt,
    arc=2pt,
    left=6pt, right=6pt, top=4pt, bottom=4pt,
}

\newtcolorbox{claudebox}{
    colback=claudebg,
    colframe=claudepurple,
    coltitle=white,
    fonttitle=\bfseries,
    title=Text added by Claude,
    boxrule=0.5pt,
    arc=2pt,
    left=6pt, right=6pt, top=4pt, bottom=4pt,
}

% --- Title ---
\title{
    \vspace{-1cm}
    \textbf{\Huge X-CAVATE}\\[8pt]
    \LARGE CLI \& Python API Reference\\[6pt]
    \large Version 2.0.0
}
\author{}
\date{\today}

\begin{document}

% --- Title page ---
\begin{titlepage}
\centering
\vspace*{3cm}

{\Huge \textbf{X-CAVATE}}\\[12pt]
{\LARGE CLI \& Python API Reference}\\[20pt]
{\large Version 2.0.0}\\[40pt]

{\large Convert vascular network geometries into collision-free\\[4pt]
3D printer toolhead pathways and G-code.}

\vspace{2cm}

{\normalsize
\textbf{SM} = Single-Material \quad|\quad \textbf{MM} = Multimaterial
}

\vspace{1cm}

{\small This document combines content from the original X-CAVATE Supplement
User Manual with new CLI and Python API documentation for the refactored
package. Sections containing new text are marked with a purple
\textcolor{claudepurple}{\textbf{``Text added by Claude''}} callout box.}

\vfill

{\normalsize \today}

\end{titlepage}

% --- Table of Contents ---
\tableofcontents
\newpage

% ===================================================================
\section{Introduction}
\label{sec:intro}
% ===================================================================

\begin{claudebox}
X-CAVATE (Cross-Collision Avoidance for VAscular Tree Embedding) converts
vascular network geometries exported from SimVascular into collision-free 3D
printer toolhead pathways and G-code. The software supports three usage modes:

\begin{enumerate}[leftmargin=*]
    \item \textbf{Streamlit GUI} --- A web-based graphical interface requiring
    no coding experience. See the companion \textit{GUI User Guide} for full
    documentation.
    \item \textbf{Command-Line Interface (CLI)} --- The \code{xcavate} command
    for scripted and batch workflows. Documented in
    Section~\ref{sec:cli}.
    \item \textbf{Python API} --- Import \code{xcavate} as a library for
    programmatic access. Documented in Section~\ref{sec:python-api}.
\end{enumerate}

This document serves as the reference manual for the CLI and Python API. It
also preserves the theoretical background and workflow descriptions from the
original supplement user manual.
\end{claudebox}

% ===================================================================
\section{Setup}
\label{sec:setup}
% ===================================================================

\subsection{Original Setup Instructions}

% PRESERVED from supplement lines 22-35 (updated URL)
Navigate to \url{https://github.com/sohams-MASS/X-CAVATE}.
Click on the green ``Code'' button and download the \code{.zip} file.
Place the downloaded folder in the directory of your choice, such as the
Desktop.

[See original supplement for screenshot.]

\begin{claudebox}
The repository URL has changed from the original
\code{github.com/jessica-herrmann/xcavate} to the current
\code{github.com/sohams-MASS/X-CAVATE}. The folder structure has been
refactored from a single-script layout into a proper Python package.
\end{claudebox}

\subsection{Installation}

\subsubsection{Option A: Conda (Recommended)}

\begin{enumerate}[leftmargin=*]
    \item Open a terminal (Mac: Terminal.app; Windows: Anaconda Prompt or PowerShell).
    \item Confirm conda is installed:
    \begin{lstlisting}
conda --version
    \end{lstlisting}
    If not installed, download Miniconda from \url{https://docs.conda.io/en/latest/miniconda.html}.
    \item Clone the repository and enter the directory:
    \begin{lstlisting}
git clone https://github.com/sohams-MASS/X-CAVATE.git
cd X-CAVATE
    \end{lstlisting}
    \item Create and activate the conda environment:
    \begin{lstlisting}
conda env create -f environment.yml
conda activate xcavate-gui
    \end{lstlisting}
    \item Verify the installation:
    \begin{lstlisting}
xcavate --help
    \end{lstlisting}
\end{enumerate}

\subsubsection{Option B: pip}

\begin{enumerate}[leftmargin=*]
    \item Confirm Python 3.9+ is installed: \code{python3 -{}-version}
    \item Clone the repository (same as above).
    \item Install with GUI support:
    \begin{lstlisting}
pip install ".[gui]"
    \end{lstlisting}
    For CLI-only usage (no web interface): \code{pip install .}
\end{enumerate}

\subsubsection{Option C: Docker}

\begin{enumerate}[leftmargin=*]
    \item Install Docker Desktop from \url{https://www.docker.com/products/docker-desktop}.
    \item Build the image:
    \begin{lstlisting}
docker build -t xcavate .
    \end{lstlisting}
    \item Run the container:
    \begin{lstlisting}
docker run -p 8501:8501 xcavate
    \end{lstlisting}
    \item Open \url{http://localhost:8501} in your browser.
\end{enumerate}

% ===================================================================
\section{Input Files}
\label{sec:inputs}
% ===================================================================

\subsection{Inputs Overview}

% PRESERVED from supplement lines 36-56
The \textit{inputs} folder contains two subfolders: \textit{custom} and
\textit{extension}.

[See original supplement for screenshot.]

\subsubsection{Extension Folder}

For a description of how to utilize the \textit{inputs/extension} folder to
implement segment extension (gap closure), see Section~\ref{sec:gap-closure}
of this guide.

\subsubsection{Custom G-code Folder}

% PRESERVED from supplement lines 44-51
The \textit{inputs/custom} folder contains sample \code{.txt} files which can
be updated to contain custom G-code wrappers to adapt X-CAVATE to your
specific printer. For example, to provide your own G-code for header
information, replace the placeholder text within \code{header\_code.txt} with
your own code to accomplish this task.

[See original supplement for screenshot.]

The screenshot below of an output file, \code{gcode\_SM\_positiveInk.txt},
provides an example of integrating custom G-code into the output G-code
\code{.txt} file generated for single material printing with a positive ink
displacement-based printer. Editing the \code{header\_code.txt} file in
advance of running X-CAVATE will replace the placeholder text in the first
line with the custom G-code saved within \code{header\_code.txt}.

[See original supplement for screenshot.]

Place your two \code{.txt} output files from SimVascular, one containing the
network splines as coordinate lists, and the other containing the identities of
the inlet and outlet nodes, into the \textit{inputs} folder.

[See original supplement for screenshot.]

\subsection{Network File Format}

\begin{lstlisting}
Vessel: 0, Number of Points: 100
33.4878, 0.0000, 53.5303, 0.0200
33.4115, -0.5285, 52.7311, 0.0200
...
Vessel: 1, Number of Points: 100
15.0600, -37.3300, 11.6600, 0.0190
...
\end{lstlisting}

\begin{tabularx}{\textwidth}{l l X}
\toprule
\textbf{Column} & \textbf{Required} & \textbf{Description} \\
\midrule
1--3 & Yes & x, y, z coordinates (cm) \\
4    & No  & Vessel radius (needed for speed calculation) \\
5    & No  & Arterial/venous flag: 0 = venous, nonzero = arterial (needed for multimaterial) \\
\bottomrule
\end{tabularx}

\begin{warnbox}
Coordinates in the network file must be in \textbf{centimeters} --- X-CAVATE
converts to millimeters internally (multiplication by 10).
\end{warnbox}

\subsection{Inlet / Outlet File Format}

\begin{lstlisting}
inlet
21.97, -15.43, 50.00
11.81, 0.00, 21.34
outlet
36.05, -27.19, 50.00
50.00, 0.00, 26.54
\end{lstlisting}

Coordinates are matched to the nearest network point automatically.

% ===================================================================
\section{Running X-CAVATE via CLI}
\label{sec:cli}
% ===================================================================

\begin{claudebox}
The original workflow required typing \code{python3 xcavate.py} followed by
all parameters, and optionally using the \code{command\_line\_writer.py} helper
script to generate the command string. In the refactored package, X-CAVATE is
installed as a console entry point. The command is simply \code{xcavate}:

\begin{lstlisting}
xcavate --network_file <path> --inletoutlet_file <path> [OPTIONS]
\end{lstlisting}

The \code{command\_line\_writer.py} script is no longer needed --- all
parameters are passed directly as CLI flags.
\end{claudebox}

\subsection{Basic Invocation}

\begin{lstlisting}
xcavate \
  --network_file inputs/network.txt \
  --inletoutlet_file inputs/inlet_outlet.txt \
  --multimaterial 0 \
  --tolerance_flag 0 \
  --nozzle_diameter 0.5 \
  --container_height 10 \
  --num_decimals 3 \
  --speed_calc 0 \
  --plots 1 \
  --downsample 0 \
  --custom 0 \
  --printer_type 0
\end{lstlisting}

\subsection{Complete CLI Argument Reference}

\begin{claudebox}
The following tables document every CLI argument. They are derived directly
from the argument parser in \code{xcavate/cli.py}.
\end{claudebox}

\subsubsection{Required Arguments}

\begin{longtable}{L{3.8cm} l l L{5.5cm}}
\toprule
\textbf{Flag} & \textbf{Type} & \textbf{Default} & \textbf{Description} \\
\midrule
\endfirsthead
\toprule
\textbf{Flag} & \textbf{Type} & \textbf{Default} & \textbf{Description} \\
\midrule
\endhead
\code{-{}-network\_file} & str & --- & Path to \code{.txt} file containing network coordinates. \\
\code{-{}-inletoutlet\_file} & str & --- & Path to \code{.txt} file containing inlet and outlet coordinates. \\
\code{-{}-multimaterial} & int & --- & Multimaterial mode (1=yes, 0=no). \\
\code{-{}-tolerance\_flag} & float & 0 & Include tolerance? (1=yes, 0=no). \\
\code{-{}-tolerance} & float & 0 & Amount of tolerance (0 is none). \\
\code{-{}-nozzle\_diameter} & float & --- & Outer diameter of the nozzle (mm). \\
\code{-{}-container\_height} & float & --- & Height of the print container (mm). \\
\code{-{}-num\_decimals} & int & --- & Number of decimal places for rounding output values. \\
\code{-{}-speed\_calc} & int & --- & Compute print speeds for changing radii? (1=yes, 0=no). \\
\code{-{}-plots} & int & --- & Generate plots? (1=yes, 0=no). \\
\code{-{}-downsample} & int & --- & Downsample network? (1=yes, 0=no). \\
\code{-{}-custom} & int & --- & Providing custom G-code? (1=yes, 0=no). \\
\code{-{}-printer\_type} & int & 0 & Type of printer: 0=Pressure, 1=Positive Ink, 2=Aerotech. \\
\code{-{}-amount\_up} & float & 10 & Z-raise above container for nozzle switches (mm). \\
\bottomrule
\end{longtable}

\subsubsection{Geometry Parameters}

\begin{tabularx}{\textwidth}{L{3.8cm} l l X}
\toprule
\textbf{Flag} & \textbf{Type} & \textbf{Default} & \textbf{Description} \\
\midrule
\code{-{}-container\_x} & float & 50 & Container x-dimension (mm). \\
\code{-{}-container\_y} & float & 50 & Container y-dimension (mm). \\
\code{-{}-scale\_factor} & float & 1 & Network scale factor. \\
\code{-{}-top\_padding} & float & 0 & Padding above network (mm). \\
\code{-{}-downsample\_factor} & int & 1 & Downsample factor. \\
\bottomrule
\end{tabularx}

\subsubsection{Print Settings}

\begin{tabularx}{\textwidth}{L{3.8cm} l l X}
\toprule
\textbf{Flag} & \textbf{Type} & \textbf{Default} & \textbf{Description} \\
\midrule
\code{-{}-flow} & float & 0.161 & Volumetric flow rate (mm\textsuperscript{3}/s). \\
\code{-{}-print\_speed} & float & 1 & Print speed (mm/s). \\
\code{-{}-jog\_speed} & float & 5 & Jog speed (mm/s). \\
\code{-{}-jog\_speed\_lift} & float & 0.25 & Z-lift jog speed (mm/s). \\
\code{-{}-initial\_lift} & float & 0.5 & Initial lift distance (mm). \\
\code{-{}-jog\_translation} & float & 10 & Jog speed between nozzles (mm/s). \\
\code{-{}-dwell\_start} & float & 0.08 & Dwell at start (s). Aerotech only. \\
\code{-{}-dwell\_end} & float & 0.08 & Dwell at end (s). Aerotech only. \\
\bottomrule
\end{tabularx}

\subsubsection{Multimaterial Settings}

\begin{tabularx}{\textwidth}{L{3.8cm} l l X}
\toprule
\textbf{Flag} & \textbf{Type} & \textbf{Default} & \textbf{Description} \\
\midrule
\code{-{}-offset\_x} & float & 103 & Printhead x-offset (mm). \\
\code{-{}-offset\_y} & float & 0.5 & Printhead y-offset (mm). \\
\code{-{}-front\_nozzle} & int & 1 & 1 if venous nozzle in front, 2 if behind. \\
\code{-{}-printhead\_1} & str & Aa & Arterial printhead name. \\
\code{-{}-printhead\_2} & str & Ab & Venous printhead name. \\
\code{-{}-axis\_1} & str & A & Arterial z-axis name. \\
\code{-{}-axis\_2} & str & B & Venous z-axis name. \\
\code{-{}-resting\_pressure} & float & 0 & Resting pressure (psi). \\
\code{-{}-active\_pressure} & float & 5 & Active pressure (psi). \\
\bottomrule
\end{tabularx}

\subsubsection{Gap Closure}

\begin{tabularx}{\textwidth}{L{3.8cm} l l X}
\toprule
\textbf{Flag} & \textbf{Type} & \textbf{Default} & \textbf{Description} \\
\midrule
\code{-{}-num\_overlap} & int & 0 & Overlap nodes at pass boundaries. \\
\code{-{}-overlap\_algorithm} & str & retrace & Overlap algorithm: \code{retrace} (original). \\
\code{-{}-close\_sm} & int & 0 & Load gap closure file for SM? (1=yes, 0=no). \\
\code{-{}-close\_mm} & int & 0 & Load gap closure file for MM? (1=yes, 0=no). \\
\bottomrule
\end{tabularx}

\subsubsection{Positive Ink Displacement}

\begin{tabularx}{\textwidth}{L{4.5cm} l l X}
\toprule
\textbf{Flag} & \textbf{Type} & \textbf{Default} & \textbf{Description} \\
\midrule
\code{-{}-positiveInk\_start} & float & 0 & Plunger displacement at pass start. \\
\code{-{}-positiveInk\_end} & float & 0 & Plunger displacement at pass end. \\
\code{-{}-positiveInk\_radii} & int & 0 & Use per-node vessel radii? (1=yes, 0=no). \\
\code{-{}-positiveInk\_diam} & float & 1 & Fixed vessel diameter (mm). \\
\code{-{}-positiveInk\_syringe\_diam} & float & 1 & Syringe barrel diameter (mm). \\
\code{-{}-positiveInk\_factor} & float & 1 & Extrusion scaling factor. \\
\code{-{}-positiveInk\_start\_arterial} & float & 0 & Arterial plunger displacement at start. \\
\code{-{}-positiveInk\_start\_venous} & float & 0 & Venous plunger displacement at start. \\
\code{-{}-positiveInk\_end\_arterial} & float & 0 & Arterial plunger displacement at end. \\
\code{-{}-positiveInk\_end\_venous} & float & 0 & Venous plunger displacement at end. \\
\bottomrule
\end{tabularx}

\subsubsection{Advanced Flags}

\begin{claudebox}
These flags were added in the refactored version and were not present in the
original \code{xcavate.py} script.
\end{claudebox}

\begin{tabularx}{\textwidth}{L{3.8cm} l l X}
\toprule
\textbf{Flag} & \textbf{Type} & \textbf{Default} & \textbf{Description} \\
\midrule
\code{-{}-algorithm} & str & dfs & Pathfinding algorithm: \code{dfs} (depth-first search). \\
\code{-{}-output\_dir} & str & outputs & Output directory path. \\
\bottomrule
\end{tabularx}

\subsection{Example Commands}

\begin{claudebox}
The examples below are adapted from the original supplement user manual,
updated for the new \code{xcavate} CLI entry point.
\end{claudebox}

\subsubsection{Minimal Single-Material Run}

\begin{lstlisting}
xcavate \
  --network_file inputs/network.txt \
  --inletoutlet_file inputs/inlet_outlet.txt \
  --multimaterial 0 --tolerance_flag 0 \
  --nozzle_diameter 0.65 --container_height 50 \
  --num_decimals 5 --amount_up 10 \
  --speed_calc 0 --plots 1 --downsample 0 \
  --custom 0 --printer_type 0
\end{lstlisting}

\subsubsection{Multimaterial with Speed Calculation and Custom G-code}

% Adapted from supplement line 114
\begin{lstlisting}
xcavate \
  --network_file inputs/multimaterial_network.txt \
  --inletoutlet_file inputs/inlet_outlet_multimaterial.txt \
  --multimaterial 1 --tolerance_flag 0 --tolerance 0 \
  --nozzle_diameter 0.65 --num_decimals 5 \
  --amount_up 10 --container_height 50 \
  --speed_calc 1 --plots 1 --downsample 0 \
  --flow 0.127 --custom 1 --printer_type 0
\end{lstlisting}

\subsubsection{Using a Custom Output Directory}

\begin{lstlisting}
xcavate \
  --network_file inputs/network.txt \
  --inletoutlet_file inputs/inlet_outlet.txt \
  --multimaterial 0 --tolerance_flag 0 \
  --nozzle_diameter 0.5 --container_height 10 \
  --num_decimals 3 --speed_calc 0 --plots 1 \
  --downsample 0 --custom 0 --printer_type 0 \
  --output_dir my_results
\end{lstlisting}

% ===================================================================
\section{Running X-CAVATE via Python API}
\label{sec:python-api}
% ===================================================================

\begin{claudebox}
This entire section documents the Python API introduced in the refactored
package. It enables programmatic access to X-CAVATE for scripting, batch
processing, and integration into larger workflows.
\end{claudebox}

\subsection{Quick Start}

\begin{lstlisting}[style=python]
from pathlib import Path
from xcavate.config import XcavateConfig
from xcavate.pipeline import run_xcavate

config = XcavateConfig(
    network_file=Path("inputs/network.txt"),
    inletoutlet_file=Path("inputs/inlet_outlet.txt"),
    nozzle_diameter=0.5,
    container_height=10.0,
    num_decimals=3,
)

results = run_xcavate(config)

# Access results
print(f"SM passes: {len(results['print_passes_sm'])}")
print(f"Output in: {config.output_dir}/")
\end{lstlisting}

\subsection{XcavateConfig Dataclass}

The \code{XcavateConfig} dataclass (defined in \code{xcavate/config.py})
replaces the 30+ global variables from the original script. All fields with
defaults are optional.

\begin{longtable}{L{4cm} l l L{5cm}}
\toprule
\textbf{Field} & \textbf{Type} & \textbf{Default} & \textbf{Description} \\
\midrule
\endfirsthead
\toprule
\textbf{Field} & \textbf{Type} & \textbf{Default} & \textbf{Description} \\
\midrule
\endhead
\multicolumn{4}{l}{\textit{Required parameters}} \\
\midrule
\code{network\_file} & Path & --- & Path to network coordinates file. \\
\code{inletoutlet\_file} & Path & --- & Path to inlet/outlet file. \\
\code{nozzle\_diameter} & float & --- & Outer diameter of nozzle (mm). \\
\code{container\_height} & float & --- & Height of print container (mm). \\
\code{num\_decimals} & int & --- & Decimal places for output rounding. \\
\midrule
\multicolumn{4}{l}{\textit{Feature flags}} \\
\midrule
\code{amount\_up} & float & 10.0 & Z-raise between passes (mm). \\
\code{multimaterial} & bool & False & Enable multimaterial mode. \\
\code{tolerance\_flag} & bool & False & Enable tolerance. \\
\code{speed\_calc} & bool & False & Compute per-node speeds from radii. \\
\code{generate\_plots} & bool & True & Create 3D HTML visualizations. \\
\code{downsample} & bool & False & Enable downsampling. \\
\code{custom\_gcode} & bool & False & Use custom G-code templates. \\
\code{printer\_type} & PrinterType & PRESSURE & Printer type enum. \\
\code{algorithm} & PathfindingAlgorithm & DFS & Pathfinding strategy enum. \\
\midrule
\multicolumn{4}{l}{\textit{Geometry}} \\
\midrule
\code{tolerance} & float & 0.0 & Tolerance distance (mm). \\
\code{scale\_factor} & float & 1.0 & Network scale multiplier. \\
\code{top\_padding} & float & 0.0 & Padding above network top (mm). \\
\code{container\_x} & float & 50.0 & Container width (mm). \\
\code{container\_y} & float & 50.0 & Container depth (mm). \\
\code{convert\_factor} & float & 10.0 & Unit conversion (cm $\to$ mm). \\
\code{downsample\_factor} & int & 1 & Downsample factor. \\
\midrule
\multicolumn{4}{l}{\textit{Print settings}} \\
\midrule
\code{flow} & float & 0.161 & Volumetric flow rate (mm\textsuperscript{3}/s). \\
\code{print\_speed} & float & 1.0 & Print speed (mm/s). \\
\code{jog\_speed} & float & 5.0 & Jog speed (mm/s). \\
\code{jog\_speed\_lift} & float & 0.25 & Z-lift jog speed (mm/s). \\
\code{jog\_translation} & float & 10.0 & Inter-nozzle jog speed (mm/s). \\
\code{initial\_lift} & float & 0.5 & Initial lift distance (mm). \\
\code{dwell\_start} & float & 0.08 & Dwell at start (s). \\
\code{dwell\_end} & float & 0.08 & Dwell at end (s). \\
\midrule
\multicolumn{4}{l}{\textit{Multimaterial}} \\
\midrule
\code{offset\_x} & float & 103.0 & Printhead x-offset (mm). \\
\code{offset\_y} & float & 0.5 & Printhead y-offset (mm). \\
\code{front\_nozzle} & int & 1 & 1=venous in front, 2=behind. \\
\code{printhead\_1} & str & ``Aa'' & Arterial printhead name. \\
\code{printhead\_2} & str & ``Ab'' & Venous printhead name. \\
\code{axis\_1} & str & ``A'' & Arterial z-axis name. \\
\code{axis\_2} & str & ``B'' & Venous z-axis name. \\
\code{resting\_pressure} & float & 0.0 & Resting pressure (psi). \\
\code{active\_pressure} & float & 5.0 & Active pressure (psi). \\
\midrule
\multicolumn{4}{l}{\textit{Gap closure}} \\
\midrule
\code{num\_overlap} & int & 0 & Overlap nodes at boundaries. \\
\code{overlap\_algorithm} & OverlapAlgorithm & RETRACE & Overlap algorithm enum. \\
\code{close\_sm} & bool & False & Load SM gap closure file. \\
\code{close\_mm} & bool & False & Load MM gap closure file. \\
\midrule
\multicolumn{4}{l}{\textit{Positive ink displacement}} \\
\midrule
\code{positive\_ink\_start} & float & 0.0 & Plunger displacement at pass start. \\
\code{positive\_ink\_end} & float & 0.0 & Plunger displacement at pass end. \\
\code{positive\_ink\_radii} & bool & False & Use per-node vessel radii. \\
\code{positive\_ink\_diam} & float & 1.0 & Fixed vessel diameter (mm). \\
\code{positive\_ink\_syringe\_diam} & float & 1.0 & Syringe barrel diameter (mm). \\
\code{positive\_ink\_factor} & float & 1.0 & Extrusion scaling factor. \\
\code{positive\_ink\_start\_arterial} & float & 0.0 & Arterial start displacement. \\
\code{positive\_ink\_start\_venous} & float & 0.0 & Venous start displacement. \\
\code{positive\_ink\_end\_arterial} & float & 0.0 & Arterial end displacement. \\
\code{positive\_ink\_end\_venous} & float & 0.0 & Venous end displacement. \\
\midrule
\multicolumn{4}{l}{\textit{Paths}} \\
\midrule
\code{custom\_gcode\_dir} & Path & inputs/custom & Directory for custom G-code templates. \\
\code{extension\_dir} & Path & inputs/extension & Directory for gap extension files. \\
\code{output\_dir} & Path & outputs & Output directory. \\
\bottomrule
\end{longtable}

\subsection{Enums}

\begin{lstlisting}[style=python]
from xcavate.config import PrinterType, PathfindingAlgorithm, OverlapAlgorithm

# PrinterType
PrinterType.PRESSURE       # 0 - Pneumatic extrusion
PrinterType.POSITIVE_INK   # 1 - Positive ink displacement
PrinterType.AEROTECH       # 2 - Aerotech 6-axis controller

# PathfindingAlgorithm
PathfindingAlgorithm.DFS        # "dfs" - Modified depth-first search
# Only DFS is currently supported

# OverlapAlgorithm
OverlapAlgorithm.RETRACE      # "retrace" - Scan all previous passes
\end{lstlisting}

\subsection{run\_xcavate() Function}

\begin{lstlisting}[style=python]
def run_xcavate(
    config: XcavateConfig,
    progress_cb: Optional[Callable[[int, str], None]] = None,
) -> dict:
\end{lstlisting}

\textbf{Parameters:}
\begin{description}[leftmargin=2cm]
    \item[\code{config}] Fully populated \code{XcavateConfig} instance.
    \item[\code{progress\_cb}] Optional callback \code{(step\_number, description)}
    for progress reporting. Steps are numbered 1--7.
\end{description}

\textbf{Returns} a dict with the following keys:

\begin{tabularx}{\textwidth}{l X}
\toprule
\textbf{Key} & \textbf{Description} \\
\midrule
\code{print\_passes\_sm} & Single-material print passes (dict of pass index $\to$ node list). \\
\code{print\_passes\_mm} & Multimaterial print passes (or \code{None}). \\
\code{points} & Interpolated coordinate array (numpy ndarray). \\
\code{points\_original} & Original coordinate array before interpolation. \\
\code{coord\_num\_dict\_original} & Original vessel $\to$ point count mapping. \\
\code{changelog} & List of changelog lines from gap closure. \\
\code{speed\_map\_sm} & Speed map for SM (or \code{None}). \\
\code{speed\_map\_mm} & Speed map for MM (or \code{None}). \\
\code{material\_map} & Material classification per pass (or \code{None}). \\
\code{instructions} & Operator instructions dict (start position, padding, etc.). \\
\bottomrule
\end{tabularx}

\subsection{Complete Example}

\begin{lstlisting}[style=python]
from pathlib import Path
from xcavate.config import (
    XcavateConfig,
    PrinterType,
    PathfindingAlgorithm,
)
from xcavate.pipeline import run_xcavate

config = XcavateConfig(
    network_file=Path("inputs/multimaterial_network.txt"),
    inletoutlet_file=Path("inputs/inlet_outlet_mm.txt"),
    nozzle_diameter=0.65,
    container_height=50.0,
    num_decimals=5,
    multimaterial=True,
    speed_calc=True,
    flow=0.127,
    generate_plots=True,
    custom_gcode=True,
    printer_type=PrinterType.PRESSURE,
    algorithm=PathfindingAlgorithm.DFS,
    output_dir=Path("my_output"),
)

results = run_xcavate(config)

# Inspect results
sm = results["print_passes_sm"]
mm = results["print_passes_mm"]
print(f"Generated {len(sm)} SM passes")
if mm:
    print(f"Generated {len(mm)} MM passes")

# Access start position for operator instructions
instr = results["instructions"]
pos = instr["start_position"]
print(f"Start at X={pos['x']}, Y={pos['y']}, Z={pos['z']}")
\end{lstlisting}

% ===================================================================
\section{Calibration}
\label{sec:calibration}
% ===================================================================

\subsection{Updating Vessel Radii}

% PRESERVED from supplement lines 164-178 (adapted for new workflow)
Follow the calibration protocol to create \code{diameters.csv} and
\code{diameters\_validation.csv}. In the original workflow, these files were
placed within the \textit{xcavate-main/inputs} folder and processed by running
\code{python3 calibration.py} from the command line. The computed volumetric
flow rate, $Q$, was output to the command line interface and saved within a
subfolder \textit{xcavate-main/outputs/radii}.

[See original supplement for screenshots.]

The calibration curve appears within \code{radii\_pressure.html}.

[See original supplement for screenshot.]

The percent error outputs to the command line interface, as well as saves in
tabular format to the output folder as \code{error\_pressure.html}.

[See original supplement for screenshots.]

\begin{claudebox}
In the refactored version, calibration is performed through either:
\begin{itemize}[nosep]
    \item The \gui{Calibration} tab in the Streamlit GUI (see GUI User Guide).
    \item The Python API functions described below.
\end{itemize}
The standalone \code{calibration.py} and \code{calibration\_validation.py}
scripts have been replaced by the \code{xcavate.core.calibration} module.
\end{claudebox}

\subsection{Q Computation Equations}

% PRESERVED from supplement lines 220-235
To incorporate the specified vessel radii, X-CAVATE continuously updates the
print speed of each coordinate within the network, at constant pressure, to
dynamically alter the printed filament diameter. Computation of the print
speed depends on the following definition of volumetric flow rate:

\textbf{Equation 1:}
\begin{equation}
Q = \frac{V}{t}
\end{equation}

where $Q$ is the volumetric flow rate of the ink through the nozzle orifice,
$V$ is the volume of the printed filament, and $t$ is the duration of the
print. Assuming that the cylindrical nozzle extrudes a cylindrical filament,
$V$ becomes $\pi r^2 L$, where $r$ is the filament radius and $L$ is the
length of the printed filament. Substituting into Equation~1 results in:

\textbf{Equation 2:}
\begin{equation}
Q = \frac{\pi r^2 L}{t}
\end{equation}

Since $L$ is equivalent to the nozzle translation velocity, $v$, multiplied
by the duration of travel, $t$, Equation~2 becomes:

\textbf{Equation 3:}
\begin{equation}
Q = \frac{\pi r^2 v t}{t}
\end{equation}

Rearrangement of Equation~3 results in:

\textbf{Equation 4:}
\begin{equation}
r = \sqrt{\frac{Q}{\pi v}}
\end{equation}

Thus, for a known filament radius $r$ and ink volumetric flow rate $Q$, the
nozzle translation velocity $v$ can be calculated.

\subsection{Python API: Calibration Functions}

\begin{claudebox}
The calibration functions below are defined in
\code{xcavate/core/calibration.py} and can be called directly from Python.
\end{claudebox}

\subsubsection{compute\_flow\_rate()}

\begin{lstlisting}[style=python]
from xcavate.core.calibration import compute_flow_rate

result = compute_flow_rate(
    measured_csv_bytes=open("diameters.csv", "rb").read(),
    print_speeds=[0.5, 1.0, 2.0],     # mm/s, one per group
    measurements_per_target=3,          # rows to average per group
    output_dir=Path("outputs/radii"),   # optional: save to disk
)

print(f"Q = {result['q_value']:.4f} mm^3/s")
# result also contains: 'figure', 'q_txt_bytes', 'html_bytes'
\end{lstlisting}

\textbf{Parameters:}
\begin{description}[leftmargin=2cm]
    \item[\code{measured\_csv\_bytes}] Raw bytes of a CSV file with a
    \code{Length} column containing measured filament diameters in
    \textbf{microns}.
    \item[\code{print\_speeds}] Print speeds in mm/s used during calibration,
    one per group. Pressure is held constant; speed is varied.
    \item[\code{measurements\_per\_target}] Number of rows in the CSV to
    average for each print speed group (default: 3).
    \item[\code{output\_dir}] If provided, \code{Q\_pressure.txt} and
    \code{radii\_pressure.html} are written to this directory.
\end{description}

\textbf{Returns} a dict with keys: \code{q\_value} (float), \code{figure}
(Plotly Figure), \code{q\_txt\_bytes} (bytes), \code{html\_bytes} (bytes).

\subsubsection{validate\_calibration()}

\begin{lstlisting}[style=python]
from xcavate.core.calibration import validate_calibration

result = validate_calibration(
    target_diameters=[0.5, 0.8, 1.0],   # expected diameters in mm
    measured_csv_bytes=open("diameters_validation.csv", "rb").read(),
    measurements_per_target=3,
)

for t, m, e in zip(
    result["target"], result["measured_avg"], result["percent_error"]
):
    print(f"Target: {t} mm, Measured: {m} mm, Error: {e}%")
# result also contains: 'csv_bytes', 'html_bytes', 'figure'
\end{lstlisting}

\textbf{Parameters:}
\begin{description}[leftmargin=2cm]
    \item[\code{target\_diameters}] Expected diameters in mm.
    \item[\code{measured\_csv\_bytes}] Raw bytes of a CSV file with a
    \code{Length} column containing measured diameters in \textbf{microns}.
    \item[\code{measurements\_per\_target}] Number of rows to average per
    target diameter (default: 3).
\end{description}

\textbf{Returns} a dict with keys: \code{target}, \code{measured\_avg},
\code{percent\_error} (lists), \code{csv\_bytes} (bytes),
\code{html\_bytes} (bytes), \code{figure} (Plotly Figure).

% ===================================================================
\section{Outputs}
\label{sec:outputs}
% ===================================================================

\subsection{Output Overview}

% PRESERVED from supplement lines 100-102
The output from X-CAVATE automatically saves within the \textit{outputs}
folder (or the directory specified by \code{-{}-output\_dir}), which, prior to running X-CAVATE,
is empty. Output takes the form of G-code files, network plots, the
mathematical graph describing the network, and a changelog. The plots depict
the individual print passes (if the plots parameter was enabled), and the
\textit{graph} folder contains individual files breaking down the network into
various formats, including nodes, edges, and coordinates, for ease of future
data manipulation.

[See original supplement for screenshot.]

\subsection{Directory Structure}

\begin{lstlisting}
outputs/
  graph/
    graph.txt               # Adjacency graph
    special_nodes.txt       # Branchpoints, endpoints, inlets, outlets
    SM_x_coords.txt         # Single-material x coordinates per pass
    SM_y_coords.txt         # Single-material y coordinates per pass
    SM_z_coords.txt         # Single-material z coordinates per pass
    SM_combined.txt         # Combined xyz per pass
    MM_x_coords.txt         # Multimaterial x coordinates (if enabled)
    MM_y_coords.txt         # Multimaterial y coordinates (if enabled)
    MM_z_coords.txt         # Multimaterial z coordinates (if enabled)
    MM_combined.txt         # Multimaterial combined xyz (if enabled)
  gcode/
    gcode_SM_pressure.txt   # Single-material G-code
    gcode_MM_pressure.txt   # Multimaterial G-code (if enabled)
  plots/
    network_original.html   # 3D plot of original network
    network_SM.html         # 3D plot of SM print passes
    network_MM.html         # 3D plot of MM print passes (if enabled)
  changelog.txt             # Gap closure operations log
\end{lstlisting}

\begin{infobox}
The G-code filename suffix matches your selected printer type:
\code{pressure}, \code{positive\_ink}, or \code{aerotech}.
\end{infobox}

\subsection{Output File Table}

% PRESERVED from supplement lines 125-148
\begin{longtable}{L{4cm} L{2cm} L{8cm}}
\toprule
\textbf{File} & \textbf{Directory} & \textbf{Description} \\
\midrule
\endfirsthead
\toprule
\textbf{File} & \textbf{Directory} & \textbf{Description} \\
\midrule
\endhead
\code{graph.txt} & \code{graph/} & List of nodes and edges constituting the mathematical graph describing the network. \\
\code{special\_nodes.txt} & \code{graph/} & Node numbers corresponding to ``special nodes'' within the graph, i.e.\ branchpoints, endpoints, inlets, outlets, and daughter nodes. \\
\code{SM\_x\_coords.txt} & \code{graph/} & Cartesian $x$-coordinates corresponding to nodes within the network, grouped by print pass and print order (for single material printing). \\
\code{SM\_y\_coords.txt} & \code{graph/} & Cartesian $y$-coordinates, grouped by print pass and print order (for single material printing). \\
\code{SM\_z\_coords.txt} & \code{graph/} & Cartesian $z$-coordinates, grouped by print pass and print order (for single material printing). \\
\code{SM\_combined.txt} & \code{graph/} & List of Cartesian coordinates describing the network, grouped by individual print pass for the passes constituting the single material print. \\
\code{MM\_*.txt} & \code{graph/} & Multimaterial equivalents of the above (if enabled). \\
\code{gcode\_SM\_*.txt} & \code{gcode/} & Single-material G-code for the selected printer type. \\
\code{gcode\_MM\_*.txt} & \code{gcode/} & Multimaterial G-code (if enabled). \\
\code{network\_original.html} & \code{plots/} & Interactive 3D visualization of the original vascular network. \\
\code{network\_SM.html} & \code{plots/} & Interactive 3D visualization of SM print passes. \\
\code{network\_MM.html} & \code{plots/} & Interactive 3D visualization of MM print passes. \\
\code{changelog.txt} & root & Log of all gap closure operations performed. \\
\bottomrule
\end{longtable}

\subsection{Changelog}

% PRESERVED from supplement lines 105-107
The \code{changelog.txt} file logs progress through X-CAVATE, allowing the
user to understand exactly how the mathematical graph defining the vascular
network has been created and manipulated by the program. For a description of
the various pre-processing and post-processing steps involved in manipulating
the network, see Section~\ref{sec:network-processing}.

[See original supplement for screenshot.]

\subsection{Plots}

% PRESERVED from supplement lines 150-162
The \textit{plots} folder contains visual plots of the networks, subdivided by
print passes.

[See original supplement for screenshots.]

The plots enable the user to step through the print passes chronologically,
visually replicating the order in which the printer will construct the network.

[See original supplement for screenshot of stepping through individual print passes.]

% ===================================================================
\section{Gap Closure / Segment Extension}
\label{sec:gap-closure}
% ===================================================================

% PRESERVED from supplement lines 181-190
\subsection{Identifying Print Passes for Segment Extension}

To identify an individual print pass for segment extension, open the network
plot after running X-CAVATE. Hover over the endpoint node of interest; an
informational box will appear which contains the x, y, z coordinates of the
node, along with the print path (e.g., trace~35 refers to print pass~35).

[See original supplement for screenshot.]

To create additional extensions for single material printing, replace the
placeholder text within \code{pass\_to\_extend\_SM.txt} and
\code{deltas\_to\_extend\_SM.txt}. Use the \code{pass\_to\_extend\_SM.txt}
file to identify the numbers of the print passes, one per line, containing the
endpoints to extend. Use the \code{deltas\_to\_extend\_SM.txt} to specify the
linear distance by which to extend the $x$-, $y$-, and $z$-coordinates of the
endpoint. In the example below, print pass~4 will be extended by 2.0~mm in the
$x$-direction, 3.5~mm in the $y$-direction, and 1.4~mm in the $z$-direction.
Print pass~6 will be extended by 5.2~mm in the $x$-direction, 1.3~mm in the
$y$-direction, and 2.4~mm in the $z$-direction. The process is the same for the
multimaterial (\_MM) files.

[See original supplement for screenshots.]

% ===================================================================
\section{X-CAVATE Structure: Description of Network Processing}
\label{sec:network-processing}
% ===================================================================

% PRESERVED verbatim from supplement lines 192-241

\subsection{Pre-processing}

X-CAVATE accepts an input vascular network in the form of two text (\code{.txt}) files.
One of these files specifies the network in terms of the Cartesian coordinates
composing its vessel splines, and the other specifies the coordinates of the
inlets and outlets. The first text file, named ``VascularNetwork.txt,''
consists of a minimum of 3 columns and a maximum of 5 columns. Column~1
contains the $x$-coordinates, column~2 contains the $y$-coordinates, and
column~3 contains the $z$-coordinates; these three columns are mandatory.
Column~4, if included, contains the radii measurements. Column~5, if included,
contains the assignment as either arterial (represented with a ``0'') or
venous (represented with a ``1''). Each row of VascularNetwork.txt therefore
specifies an individual coordinate within one of the splines constituting the
vascular network model. The coordinates are further grouped within the text
file by vessel number, which is specified as a header above each sublist of
coordinates. The second text file, called ``InletOutlet.txt,'' specifies which
coordinates within the network are the network inlets and outlets. The
coordinates within InletOutlet.txt must match their appearance within
VascularNetwork.txt exactly, meaning, for example, that if an inlet coordinate
contains 10 digits within InletOutlet.txt, it must contain the same 10 digits
within VascularNetwork.txt.

Upon importation of the network, X-CAVATE reformats the input file content into
a graph. In a graph, coordinates termed ``nodes'' are connected to one another
by lines, called ``edges.'' Any two nodes which are connected to one another by
one edge are termed ``neighbors.'' Each node is also its own neighbor. Edges
may have an orientation, or defined direction of traversal between its two
nodes, in which case the graph is called ``directed.'' Graphs whose edges are
unoriented, or directionless, are called ``undirected,'' and graphs consisting
of a combination of oriented and unoriented edges are called ``mixed'' (Tarjan).

To label each coordinate with a nodal identity, the coordinates within the
network are consecutively assigned whole numbers beginning with ``0'' until all
nodes within the network have been assigned a number. To assist with parsing the
network into subgroups of neighboring nodes, we then introduce specific
terminology to define the types of nodes within the network. These include
``endpoint'' nodes, ``parent branch point'' nodes, ``daughter'' nodes, and
``standard'' nodes. ``Endpoint'' nodes are nodes which are either
the first node or the last node of an individual vessel. Inlet and outlet nodes
are both forms of endpoint nodes which are only connected by one edge to one
node within the network. ``Parent branch point'' nodes are endpoint nodes which
form the first node of a daughter vessel which branches off of an existing
parent vessel. The parent branch point node always exists on the parent vessel
from which it branches, meaning that it exists on the line formed between its
two ``daughter'' nodes. The ``daughter'' nodes are members of the parent vessel
and are neighbors of the parent branch point. Parent branch points therefore
always have four neighbors: the two ``daughter'' nodes, itself, and the node to
which it is connected along the daughter vessel (i.e., the second node within
the daughter vessel). This latter node is a ``standard'' node, which is the
term given to any node within the network that is not an endpoint, a parent
branch point, or a daughter node. The printed ink filaments are represented by
the edges within the graph.

Prior to generating the graph, X-CAVATE performs a linear interpolation to
ensure that no adjacent coordinates within individual vessels are separated by
a distance greater than the nozzle radius. This step introduces additional
coordinates to allow the nodes to approximate the edges, which assists with
collision detection during the modified depth-first search portion of the
algorithm. The linear interpolation does not affect the network geometry
because, in each case, it introduces the coordinates along the line which
already exists between the original adjacent nodes.

\subsection{Modified Depth-First Search (DFS)}

Following interpolation and structuring of the graph, X-CAVATE performs a
modified depth-first search (DFS). DFS is a standard algorithm for traversing
all nodes within a graph. Initially, all nodes within the graph are marked as
``unvisited,'' meaning that they have not been traversed by the algorithm. To
explore the graph, DFS begins from an initial node and traverses one of its
connected edges to a neighboring node. From that node, DFS selects an
untraveled edge and repeats this process until eventually reaching a node whose
neighbors have all been visited. When this occurs, DFS backtracks along the
traversed edges emanating from that node until it encounters a node with an
unvisited neighbor. It then proceeds along that edge and continues in this
manner until all nodes within the graph have been visited. DFS can run
iteratively or recursively (Tarjan).

X-CAVATE adapts DFS to identify collision-free traversal of vascular networks.
The first modification is the introduction of a validity function, which
determines whether, on a given DFS iteration, a node is ``invalid.'' An
invalid node within this network is one that, if printed, would result in an
eventual collision between itself and the nozzle when the nozzle visits a
future, as-yet unvisited coordinate. To identify invalid coordinates,
X-CAVATE constructs a circular boundary region around the node which matches
the dimensions of the nozzle cross-section. It then extends this region
infinitely downward in the $z$-direction to form a cylindrical region. If any
unvisited nodes within the network fall within this cylindrical boundary, the
current node is marked ``invalid.'' The encountering of an invalid node
concludes the traversal pathway of the current DFS iteration, completing the
``print pass,'' and triggers a new iteration. This leads to the second DFS
modification: the new iteration will always begin at the positionally lowest
unvisited node within the network. This imposed directionality minimizes
collisions by ensuring that print pass initiation proceeds from bottom-to-top
within the network. While the individual coordinates within the print pass
are often entirely printed bottom-to-top, some passes contain nozzle
traversals that both ascend and descend in the $z$-direction. X-CAVATE
incorporates bottom-to-top directionality at the DFS level, rather than at the
edge level, to enable the printing of these segments from top-to-bottom when
permissible (i.e., when doing so does not manifest in collisions).

Each iteration of the DFS becomes an individual ``print pass,'' or segment of
a vessel which will be printed continuously. The order of nodes within an
individual print pass specifies the order in which the nozzle will traverse
them. G-code will dictate pressure-based ink extrusion between the nodes to
produce the print pass filament. Following completion of a single print pass,
G-code and Aerobasic commands will pause the pressure application, lift the
nozzle vertically out of the support bath, translate the nozzle in the
$xy$-plane to the starting ($X$, $Y$) coordinates of the next print pass,
lower the nozzle to the starting node, and reapply pressure to begin extrusion
for printing the subsequent print pass.

\subsection{Post-processing}

X-CAVATE incorporates post-processing steps to handle two features of the DFS
algorithmic output. First, because the DFS algorithm incorporates
backtracking, it occasionally traverses non-neighboring nodes, i.e., pathways
external to the intended vascular design. So, following the DFS, X-CAVATE
subdivides the output print passes into additional print passes at these
specific instances. Second, although the DFS ensures traversal of all
nodes within the network, it does not guarantee traversal of all edges.
Non-traversed edges manifest in physical gaps within the printed ink filaments.
Gaps, when they arise, occur only between the endpoint of an individual print
pass and its neighbor(s). For this reason, X-CAVATE includes post-processing
steps to sequentially check the first and last nodes of each print pass for
unconnected neighbors.

To close the gaps, X-CAVATE first searches for disconnects between standard
nodes and other nodes, excluding parent branch points. It does this by
checking whether the first or last node of the current print pass has any
neighbors within previously-printed passes which are not adjacent to it within
the current print pass (``Condition~\#0''). If so, it appends that neighbor
to the start or end of the current print pass, respectively. Second, X-CAVATE
searches for untraversed edges between daughter and parent nodes. For any
print passes with one or more daughter nodes serving as endpoints, it searches
for the corresponding parent branch point within previous print passes
(``Branchpoint Condition~\#1'') and, after appending the relevant nodes,
within future print passes (``Branchpoint Condition~\#2''). Next, X-CAVATE
identifies print passes whose endpoints consist of one or more parent branch
points (``Branchpoint Condition~\#3''). If the parent branch point does not
reside next to its non-daughter neighbor in the current print pass, X-CAVATE
searches for the non-daughter neighbor within the previous print passes and
appends as before. Finally, X-CAVATE searches for any final disconnects
remaining within the network by identifying nodes which appear only once within
the set of print passes and also appear as an endpoint (excluding inlets and
outlets). These disconnects typically take the form of gaps between
standard-node endpoints and parent branch points embedded within existing print
passes.

\subsection{Radius Computation}

% PRESERVED from supplement lines 218-235
To incorporate the specified vessel radii, X-CAVATE continuously updates the
print speed of each coordinate within the network, at constant pressure, to
dynamically alter the printed filament diameter. The computation of the print
speed depends on the volumetric flow rate equations presented in
Section~\ref{sec:calibration}. For a known filament radius $r$ and ink
volumetric flow rate $Q$, the nozzle translation velocity $v$ can be
calculated from Equation~4.

\subsection{Generation of Print Passes for Multimaterial Printing}

% PRESERVED from supplement lines 237-241
X-CAVATE enables users to specify which of the input vessels belong to the
arterial circulation and which are venous. Here, the distinction between
arterial and venous is referred to as ``vessel type.'' To translate the
programmed vessel type to the 3D printed model, X-CAVATE generates print paths
for multi-material printing with two distinct inks: arterial ink and venous ink.
Generation of multi-material print paths occurs only if the
VascularNetwork.txt file contains a fifth column specifying the vessel types.
To create the multi-material print passes, X-CAVATE further subdivides the
post-processed print passes by vessel type, similar to the subdivision step
which follows the DFS. G-code is then used to incorporate nozzle translations
between subdivided print passes to physically switch inks.

To generate the multi-material passes, X-CAVATE iterates through the print
passes which were already generated for single material printing to locate
positions within the pass at which the vessel type of adjacent nodes differs
and subdivides the passes into separate passes to permit nozzle translations
between them (for switching ink). Prior to this subdivision, however, X-CAVATE
may switch the assigned vessel type of individual nodes to avoid segmenting
into passes consisting of single nodes. The cases in which these switches occur
are as follows: 1)~in the event that the vessel type of the first node in the
pass differs from the vessel type of the subsequent nodes in that pass, 2)~in
the event that the vessel type of the last node in the pass differs from the
vessel type of the preceding nodes in that pass, 3)~in the event that the
vessel type of an individual node within a pass differs from the vessel type of
the two nodes immediately preceding and immediately following it, and 4)~in the
event that the print pass consists of only two nodes and the vessel types of
these two nodes differ; in this case, X-CAVATE will reassign the vessel type of
the first node to match that of the second node.

% ===================================================================
\section{Module Structure}
\label{sec:module-structure}
% ===================================================================

\begin{claudebox}
The original X-CAVATE codebase consisted of a single monolithic script
(\code{xcavate.py}, approximately 5,568 lines) plus standalone calibration
scripts. The refactored version is organized as a Python package with the
following structure:
\end{claudebox}

\begin{lstlisting}
xcavate/
  __init__.py              # Package init
  __main__.py              # python -m xcavate entry point
  cli.py                   # Argument parser and CLI entry point
  config.py                # XcavateConfig dataclass and enums
  pipeline.py              # run_xcavate() orchestrator

  core/
    calibration.py          # compute_flow_rate(), validate_calibration()
    gap_closure.py          # Gap closure pipeline
    graph.py                # Graph construction and special node detection
    multimaterial.py        # Material classification and subdivision
    pathfinding.py          # DFS with KD-tree collision detection
    postprocessing.py       # Subdivision, downsampling, overlap, reordering
    preprocessing.py        # Interpolation and vessel boundary detection

  io/
    reader.py               # Network and inlet/outlet file parsing
    writer.py               # Coordinate output file writing
    gcode/
      base.py               # Base G-code writer and custom code loading
      pressure.py           # Pressure printer G-code writer
      positive_ink.py       # Positive ink displacement G-code writer
      aerotech.py           # Aerotech 6-axis G-code writer

  gui/
    app.py                  # Streamlit GUI application

  spatial/
    (spatial indexing utilities)

  viz/
    plotting.py             # Plotly 3D visualization functions
\end{lstlisting}

\begin{claudebox}
Each module corresponds to a logical stage of the original monolithic script:
\begin{description}[leftmargin=1cm,nosep]
    \item[\code{cli.py}] Replaces the argparse block from \code{xcavate.py}
    lines 41--99 and the \code{command\_line\_writer.py} helper.
    \item[\code{config.py}] Replaces the 30+ global variables from
    \code{xcavate.py} lines 101--171.
    \item[\code{pipeline.py}] The orchestrator that wires all modules together,
    replacing the main execution flow of the original script.
    \item[\code{core/calibration.py}] Replaces \code{calibration.py} and
    \code{calibration\_validation.py}.
    \item[\code{core/pathfinding.py}] Replaces the DFS implementation from the
    original script.
    \item[\code{io/gcode/}] Replaces the G-code generation sections, with
    separate writers for each printer type.
\end{description}
\end{claudebox}

\end{document}
